\subsection{Temporal resolution of prescribed-flow optimization}
\label{si:thermal_resolution}
The temporal check separates the coupled discrete optimization from
the forward response to its recovered source. It uses the original
transformer mesh, the 600~s horizon and two targets selected before
optimization: target 7 and the target with the largest weighted desired
temperature exceedance, target 15. Both targets use 4, 8, 16, 32 and
64 time slabs. Geometry, prescribed velocity, frozen properties, initial
temperature and the physical upper bound of 357.3~K remain fixed.
The original-system residual threshold is $10^{-10}$; KKT and
forward-recovery thresholds are $10^{-8}$.
Sparse-direct PDAS independently checks the four-slab solutions.

Time refinement exposes a growing mode in the original advective
temperature operator. For the homogeneous equation
$C_\ell\dot y+Ky=0$, a verified generalized eigenvalue is
$-0.0664~\mathrm{s}^{-1}$. At 32 steps over 600~s, its backward-Euler
amplification modulus is 4.08. The mode's advective contribution is
negative because the prescribed finite-element velocity is only weakly
divergence-free. This property concerns the forward thermal equation;
the reduced optimization Hessian remains positive definite.

A separate pilot adds the standard half-divergence term to the
elementwise thermal transport \citep{churbanov2012convection}.
For the prescribed finite-element velocity $\beta_h$, it uses
\begin{equation}\label{eq:thermal_skew_pilot}
 c_E\beta_h\cdot\nabla y
 +\tfrac12 c_E(\nabla\cdot\beta_h)y,
 \qquad
 \nabla\cdot\beta_h=\partial_r\beta_r+\beta_r/r+\partial_z\beta_z.
\end{equation}
For constant capacity $c_E$ on an element $E$, integration by parts gives
\begin{equation}\label{eq:thermal_skew_energy}
 \int_E c_Ey\,\beta_h\cdot\nabla y
 +\tfrac12\int_E c_E(\nabla\cdot\beta_h)y^2
 =\tfrac12\int_{\partial E}c_E(\beta_h\cdot n)y^2.
\end{equation}
The integrals use the axisymmetric volume and surface measures.
Interior contributions cancel when the normal capacity flux is continuous;
otherwise their jumps remain. Degree-five quadrature verifies this
quadratic energy identity. Diffusion, streamline diffusion and all other
inputs remain unchanged. The pilot changes the thermal discretization
and has separate records from the published timing comparisons.
Its computed amplification modes decay on all seven tested forward
grids, from 4 to 256 steps. This sparse calculation tests the computed
modes and does not enumerate the complete spectrum.

Each recovered control is held constant on its original time intervals
and applied without clipping or reoptimization. Finer forward steps
align with every source discontinuity. The detailed study retains the
original and corrected optimization attempts, forward residuals, and
all successive time refinements. The final comparison uses both targets
at the coarsest and finest optimization grids. Two consecutive maximum
temperature changes must fall below 0.05~K; temperature-bound violations
are evaluated separately at every forward time level. This criterion
measures temporal sensitivity on the fixed mesh.

All ten optimizations meet the independent residual, KKT
and forward-recovery criteria. Figure~\ref{si:fig_thermal_resolution}
compares the optimized temperatures and the refined forward responses.
All four selected source histories meet the two-consecutive-grid
temperature-change criterion by 65\,536 forward steps.
The four-slab sources give maximum violations of 3.01 and
4.94~K for the nominal and demanding targets.
The corresponding 64-slab sources give 0.37 and 0.41~K.
The bounds hold at the optimization time levels.
The unchanged-source forward responses exceed them between those levels.
The comparison concerns temporal sensitivity on the existing mesh;
spatial and physical-model accuracy require separate assessment.


\begin{figure}[pos=htbp]
\centering
\includegraphics[width=\textwidth]{figures/temporal/resolution.pdf}
\caption{The separate energy-consistent transport pilot compares
discrete optimization and unchanged-source forward trajectories.
The first panel gives the root-mean-square temperature difference from the
64-slab optimization, using linear temporal reconstruction and
axisymmetric lumped-mass weights. The second panel gives bound violations
at 65\,536 forward steps. Whiskers show the maximum temperature change from
32\,768 steps, a temporal-sensitivity measure. Nominal and demanding denote
the two preselected targets. These comparisons use the same spatial mesh.}
\label{si:fig_thermal_resolution}
\end{figure}
The \href{https://github.com/teeratornk/deflation-example/tree/prescribed-temporal-data-v1/examples/temporal_resolution}
{versioned temporal-resolution example} provides both formulations,
the selected targets, all optimization and replay outcomes,
and commands for regenerating the fields and figures.
\FloatBarrier
